\chapter{Modelo conceptual}

El modelo conceptual del sistema se compone de nueve entidades: \nem{EMPRESA}, \nem{SUCURSAL}, \nem{CLIENTE}, \nem{AVAL}, \nem{INVERSION}, \nem{INTERES}, \nem{PAGO}, \nem{LIQUIDACION} y \nem{BITACORA}.

\captura{diagrama_conceptual.jpg}{Diagrama conceptual del sistema}{diagrama-conceptual}

\begin{table}[H]
	\centering
	\small
	\caption{Entidades del modelo conceptual y sus llaves}
	\label{tab:entidades}
	\begin{tabular}{@{}l l@{}}
		\toprule
		\textbf{Entidad} & \textbf{Llave} \\
		\midrule
		\nem{EMPRESA}     & PK \nem{empresa\us rfc} \\
		\nem{SUCURSAL}    & PK \nem{sucursal\us num} \\
		\nem{CLIENTE}     & PK \nem{cliente\us num}; AK \nem{cliente\us rfc} \\
		\nem{AVAL}        & PK \nem{aval\us num}; AK \nem{aval\us rfc} \\
		\nem{INVERSION}   & PK \nem{inversion\us num}; AK \nem{pagare\us num} \\
		\nem{INTERES}     & PK (\nem{inversion\us num}, \nem{interes\us fecha\us inicial}) \\
		\nem{PAGO}        & PK (\nem{inversion\us num}, \nem{interes\us fecha\us inicial}, \nem{pago\us fecha}) \\
		\nem{LIQUIDACION} & PK \nem{inversion\us num} \\
		\nem{BITACORA}    & PK (\nem{bitacora\us num}, \nem{bitacora\us fecha}) \\
		\bottomrule
	\end{tabular}
\end{table}

\clearpage

Las relaciones entre las entidades son las siguientes:

\begin{itemize}[leftmargin=*]
	\item \nem{EMPRESA} $\rightarrow$ \nem{SUCURSAL}: cada sucursal pertenece a la empresa (\nem{empresa\us rfc}).
	\item \nem{SUCURSAL} $\rightarrow$ \nem{INVERSION}: cada inversi\'on registra la sucursal donde se suscribi\'o el pagar\'e (\nem{sucursal\us num}).
	\item \nem{CLIENTE} $\rightarrow$ \nem{INVERSION}: cada inversi\'on registra al cliente que invierte (\nem{cliente\us num}).
	\item \nem{AVAL} $\rightarrow$ \nem{INVERSION}: cada inversi\'on registra, si existe, al aval que garantiza el pago (\nem{aval\us num}).
	\item \nem{INVERSION} $\rightarrow$ \nem{INTERES}: cada periodo de inter\'es pertenece a una inversi\'on (\nem{inversion\us num}).
	\item \nem{INTERES} $\rightarrow$ \nem{PAGO}: cada pago cubre un periodo de inter\'es (\nem{inversion\us num}, \nem{interes\us fecha\us inicial}).
	\item \nem{INVERSION} $\rightarrow$ \nem{LIQUIDACION}: la liquidaci\'on registra la devoluci\'on del capital de una inversi\'on (\nem{inversion\us num}).
	\item \nem{BITACORA} no se relaciona con las dem\'as entidades.
\end{itemize}

\paragraph{Reglas de negocio.} Las siguientes reglas provienen del pagaré de inversión y de la operación de la empresa, y se reflejan en el modelo:

\begin{itemize}[leftmargin=*]
	\item El n\'umero de pagar\'e es \'unico y sigue el formato consecutivo/año.
	\item El lugar de suscripci\'on del pagar\'e es la ciudad de la sucursal donde se suscribe.
	\item El aval, cuando existe, garantiza el pago de las obligaciones de la empresa y debe ser una persona distinta del cliente de la misma inversi\'on.
	\item La fecha de vencimiento es posterior a la fecha de suscripci\'on del pagar\'e.
	\item La fecha de vencimiento coincide con el día de pago, de modo que el último pago de intereses se realiza junto con la devolución del capital.
	\item Las tasas ordinaria y moratoria se registran como tasas mensuales. Cuando el pagar\'e pacta una tasa anual, se registra dividida entre doce.
	\item El inter\'es ordinario se paga mensualmente en un d\'ia fijo, entre el 1 y 28 de cada mes, a una tasa fija. Como el capital se devuelve en una sola exhibici\'on, el saldo insoluto sobre el que se calcula el inter\'es es siempre el capital invertido.
	\item El inter\'es moratorio se aplica sobre los saldos vencidos y no pagados, tanto de inter\'es como de capital, a la tasa moratoria establecida en el pagar\'e.
	\item El capital se devuelve en una sola exhibici\'on, por el monto total invertido, en la fecha de vencimiento o despu\'es de ella. Una devoluci\'on posterior al vencimiento genera inter\'es moratorio.
	\item Una inversi\'on se encuentra en uno de tres estados: \textbf{vigente}, si no tiene liquidaci\'on registrada y su fecha de vencimiento no ha pasado; \textbf{vencida}, si no tiene liquidaci\'on registrada y su fecha de vencimiento ya pas\'o; y \textbf{liquidada}, si tiene liquidaci\'on registrada.
	\item La moneda del pagar\'e es, por defecto, el peso mexicano (MXN).
\end{itemize}

\paragraph{Decisiones de diseño.} El modelo adopta las siguientes decisiones:

\begin{itemize}[leftmargin=*]
	\item \textbf{Interés programado separado del pago realizado.} \nem{INTERES} contiene lo que la empresa debe pagar en cada periodo y \nem{PAGO} lo que efectivamente pagó. La diferencia entre ambos permite obtener los estados de cuenta y determinar la mora.
	\item \textbf{Liquidación como entidad propia.} La devolución del capital se registra en \nem{LIQUIDACION}, separada de los pagos de intereses, junto con el interés moratorio que genere su retraso.
	\item \textbf{Empresa a través de la sucursal.} La inversión se relaciona con la sucursal y, por medio de ella, con la empresa, de modo que los datos de la empresa no se repiten en cada inversión.
	\item \textbf{Tasa fija en la inversión.} Como el pagaré establece una tasa fija, las tasas ordinaria y moratoria se almacenan en \nem{INVERSION} en lugar de un catálogo de tasas independiente.
	\item \textbf{Llaves sustitutas.} \nem{SUCURSAL}, \nem{CLIENTE}, \nem{AVAL} e \nem{INVERSION} usan una clave numérica como llave primaria; el RFC y el número de pagaré se conservan como llaves alternas. \nem{EMPRESA} usa su RFC como llave natural.
	\item \textbf{Clave de partición en la llave primaria.} Las llaves primarias de \nem{PAGO} y \nem{BITACORA} incluyen su fecha, porque ambos gestores exigen que la columna de partición forme parte de la llave primaria en la instancia particionada.
	\item \textbf{Bitácora independiente.} \nem{BITACORA} registra los cambios de todas las tablas sin relacionarse con ellas, de modo que sus registros se conservan aunque los datos originales se modifiquen o eliminen.
\end{itemize}